\documentclass[12pt,a4paper]{article}

%% PACHETE MATEMATICE
\usepackage{amsmath,amssymb,amsfonts}
\usepackage{amsthm}
\usepackage{mathpazo}
\usepackage{eulervm}
\usepackage{enumerate}
\usepackage{color}
\usepackage{graphicx}
\usepackage[all]{xy}
\usepackage{xcolor}
\usepackage{framed}
\usepackage[unicode]{hyperref}
\setcounter{MaxMatrixCols}{20}
\usepackage{multicol}
% \usepackage[romanian]{babel}


%% ASPECT
\linespread{1.05}
\usepackage[T1]{fontenc}
\usepackage[utf8x]{inputenc}
\usepackage[margin=2cm]{geometry}
% \usepackage[color]{showkeys}
\usepackage{anyfontsize}
\usepackage{titlesec}
\titleformat*{\section}{\large\bfseries}

\usepackage{fancyhdr}
\pagestyle{fancy}
\rhead{\small \color{gray} Notițe de Adrian Manea}
\lhead{\small \color{gray} ETTI-AM2, 2017---2018, M. Joița \& A. Niță}


\usepackage[autostyle = false, style = german]{csquotes}
\usepackage[romanian]{babel}
\newcommand\qq{\enquote}

%% COMENZILE MELE
\renewcommand*{\proofname}{Dem.:}
\newcommand\bb{\mathbb}
\newcommand\dr{\mathrm}
\newcommand\kal{\mathcal}
\newcommand\fk{\mathfrak}
\newcommand\op{\oplus}
\newcommand\ot{\otimes}
\newcommand\ds{\displaystyle}
\newcommand\id{\indent}
\newcommand\nid{\noindent}
\newcommand\seq{\subseteq}
\newcommand\sneq{\subsetneq}
\newcommand\speq{\supseteq}
\newcommand\spneq{\supsetneq}
\newcommand\rar{\rightarrow}
\newcommand\Rar{\Rightarrow}
\newcommand\xrar{\xrightarrow}
\newcommand\lar{\leftarrow}
\newcommand\Lar{\Leftarrow}
\newcommand\xlar{\xleftarrow}
\newcommand\lrar{\leftrightarrow}
\newcommand\Lrar{\Leftrightarrow}
\newcommand\ideal{\trianglelefteq}
\newcommand\ai{\text{ a.\^{i}. }}
\newcommand\sk{(\textit{Schi\c{t}\u{a}}) }
\newcommand\rhup{\rightharpoonup}
\newcommand\rhdn{\rightharpoondown}
\newcommand\lhup{\leftharpoonup}
\newcommand\lhdn{\leftharpoondown}
\newcommand\vid{\emptyset}
\newcommand\wh{\widehat}
\newcommand\ol{\overline}
\newcommand\NN{\mathbb{N}}
\newcommand\RR{\mathbb{R}}
\newcommand\ZZ{\mathbb{Z}}
\newcommand\QQ{\mathbb{Q}}
\newcommand\CC{\mathbb{C}}
\newcommand\Lap{\mathcal{L}}

%% STILURILE MELE
\newtheoremstyle{dotless}{}{}{\itshape}{}{\bfseries}{:}{ }{}

\theoremstyle{dotless}
\newtheorem{theorem}{Teorem\u{a}}
\newtheorem{proposition}{Propozi\c{t}ie}
\newtheorem{lemma}{Lem\u{a}}
\newtheorem{corollary}{Corolar}
\newtheorem{exercise}{Exerci\c{t}iu}

\newtheoremstyle{dotlessdr}{}{}{}{}{\bfseries}{:}{ }{}
\theoremstyle{dotlessdr}
\newtheorem{example}{Exemplu}
\newtheorem{definition}{Defini\c{t}ie}
\newtheorem{remark}{Observa\c{t}ie}

\newenvironment{solution}
{\begin{proof}[Soluție:\nopunct]}
  {\end{proof}}


\begin{document}
% \definecolor{labelkey}{RGB}{222,0,0}

\begin{center}
  {\large\textbf{Seminar 11 \\
      Transformarea Laplace}}
\end{center}

\vspace{2cm}

\section*{Noțiuni teoretice}

\indent\indent Transformata Laplace este tot o transformare integrală, ca și transformata
Fourier. Aplicațiile sale de interes vor fi în rezolvarea ecuațiilor și sistemelor
diferențiale. Astfel, cu ajutorul transformatei Laplace, vom putea aduce o ecuație sau
sistem diferențial la un caz algebric, iar apoi să recuperăm funcțiile inițiale aplicînd
transformata Laplace inversă.

Transformatele Laplace se aplică unor funcții speciale, care se numesc \emph{funcții original},
definite mai jos.

\begin{definition}\label{def:original}
  O funcție $ f : \RR \to \CC $ se numește \emph{original} (Laplace) dacă îndeplinește
  condițiile:
  \begin{enumerate}[(a)]
  \item $ f(t) = 0 $, pentru orice $ t < 0 $;
  \item $ f $ este continuă (eventual pe porțiuni) în intervalul $ [0, \infty) $;
  \item Există $ M > 0 $ și $ s_0 \geq 0 $ astfel încît
    \[
      |f(t)| \leq M e^{s_0t}, \forall t \geq 0.
    \]
  \end{enumerate}

  Vom nota cu $ \kal{O} $ mulțimea funcțiilor original (Laplace).
\end{definition}

Prima condiție din definiție va fi utilă pentru interpretarea fizică, în teoria
semnalelor. Într-adevăr, dacă privim $ f(t) $ ca un semnal dependent de timp, este
normal să cerem ca $ f(t) = 0 $ pentru $ t < 0 $. A doua și a treia condiție vor fi
utile pentru dezvoltarea teoriei matematice, cu ajutorul calculului integral. A treia
condiție se mai numește condiția de \emph{creștere exponențială de indice $ s_0 $}.
Ajungem acum la definiția principală.

\begin{definition}\label{def:transf-lap}
  Fie $ f \in \kal{O} $, de indice $ s_0 $ și mulțimea:
  \[
    S(s_0) = \{ s \in \CC \mid \mathrm{Re}(s) > s_0 \}.
  \]

  Funcția:
  \[
    F : S(s_0) \to \CC, \quad F(s) = \int_0^\infty f(t) e^{-st} dt
  \]
  se numește \emph{transformata Laplace} a lui $ f $ sau \emph{imaginea Laplace} a
  originalului $ f $. Vom mai folosi notația $ F = \Lap f $ sau, echivalent,
  $ \Lap f(t) = F(s) $.
\end{definition}

Conform definiției, să remarcăm că transformata Laplace are drept domeniu $ \kal{O} $,
iar imaginea este o submulțime a lui $ \CC $.

\vspace{1cm}

Principalele proprietăți ale transformatei Laplace, care ne vor ajuta în rezolvarea
de probleme, sînt redate mai jos.

\textbf{Liniaritate:}
\[
  \Lap (\alpha f(t) + \beta g(t)) = \alpha \Lap f(t) + \beta \Lap g(t).
\]

\textbf{Teorema asemănării:}
\[
  \Lap f(\alpha t) = \frac{1}{\alpha} F \Big( \frac{s}{\alpha} \Big).
\]

\textbf{Teorema deplasării:}
\[
  \Lap \big( f(t) e^{s_0 t} \big) = F(s - s_0).
\]

\textbf{Teorema întîrzierii:} Dacă $ f \in \kal{O} $ și $ \Lap f(t) = F(s) $,
atunci pentru orice $ \tau > 0 $, transformata Laplace a întîrziatei:
\[
  f_\tau(t) =
  \begin{cases}
    0, & t < \tau \\
    f(t - \tau), & t \geq \tau
  \end{cases}
\]
este $ e^{-st} F(s) $.

\textbf{Teorema derivării imaginii:}
\[
  \Lap \big( t^n f(t) \big) = (-1)^n F^{(n)} (s).
\]

\textbf{Teorema integrării originalului:} Dacă $ f \in \kal{O}, \Lap f(t) = F(s) $ și
$ g(t) = \ds\int_0^t f(\tau) d\tau $, atunci:
\[
  \Lap g(t) = \frac{1}{s} F(s).
\]

\textbf{Teorema integrării imaginii:} Fie $ \Lap f(t) = F(s) $ și $ G $ o primitivă
a lui $ F $ în semiplanul $ S(s_0) $, cu $ G(\infty) = 0 $. Atunci:
\[
  \Lap \frac{f(t)}{t} = -G(s).
\]


\newpage
%%%%%%%%%%%%%%%%%%%%%%%%%%%%%%%%%%%%%%%%%%%%%%%%%%%%%%%%%%%%%%%%%%%%%%

\section*{Tabel de transformate Laplace}

\indent\indent În tabelul de mai jos, vom considera funcțiile $ f(t) $ ca fiind funcții original,
adică nule pentru argument negativ. Echivalent, putem gîndi $ f(t) $ ca fiind, de fapt,
înmulțite cu funcția lui Heaviside:
\[
  u(t) =
  \begin{cases}
    1, & t \geq 0 \\
    0, & t < 0 
  \end{cases}.
\]

\begin{figure}[!htbp]
  \centering
  \begin{tabular}{c | c}
    $ f(t) = \Lap^{-1} (F(s)) $ & $ F(s) = \Lap(f(t)) $ \\
    \hline\hline
    $ u(t) $ & $ \dfrac{1}{s} $ \\
                                & \\
    $ u(t - \tau) $ & $ \dfrac{1}{s} e^{-\tau s} $ \\
                                & \\
    $ t $ & $ \dfrac{1}{s^2} $ \\
                                & \\
    $ t^n $ & $ \dfrac{n!}{s^{n+1}} $ \\
                                & \\
    $ t^n e^{-\alpha t} $ & $ \dfrac{n!}{(s + \alpha)^{n+1}} $ \\
                                & \\
    $ e^{-\alpha t} $ & $ \dfrac{1}{s + \alpha} $ \\
                                & \\
    $ \sin (\omega t) $ & $ \dfrac{\omega}{s^2 + \omega^2} $ \\
                                & \\
    $ \cos (\omega t) $ & $ \dfrac{s}{s^2 + \omega^2} $ \\
                                & \\
    $ \sinh(\alpha t) $ & $ \dfrac{\alpha}{s^2 - \alpha^2} $ \\
                                & \\
    $ \cosh(\alpha t) $ & $ \dfrac{s}{s^2 - \alpha^2} $ \\
                                & \\
    $ e^{-\alpha t}\sin(\omega t) $ & $ \dfrac{\omega}{(s + \alpha)^2 + \omega^2} $ \\
                                & \\
    $ e^{-\alpha t}\cos(\omega t) $ & $ \dfrac{s + \alpha}{(s + \alpha^2) + \omega^2} $ \\
                                & \\
    $ \ln t $ & $ -\dfrac{1}{s} \big( \ln s + \gamma\footnotemark \big) $
  \end{tabular}
  \caption{\emph{Transformatele Laplace uzuale}}
\end{figure}

\footnotetext{Constanta Euler-Mascheroni, $ \gamma \simeq 0,577\dots \in \RR - \QQ $}

%%%%%%%%%%%%%%%%%%%%%%%%%%%%%%%%%%%%%%%%%%%%%%%%%%%%%%%%%%%%%%%%%%%%%%
\newpage
\section*{Exerciții}

\indent\indent 1. Să se arate că funcția:
\[
  f : \RR \to \RR, \quad %
  f(t) = \begin{cases}
    t^2 + t - 3, & t \geq 0 \\
    0, & t < 0
  \end{cases}
\]
este o funcție original.

\emph{Soluție:} Evident, $ f $ este nulă pentru argument negativ, deci prima condiție este
verificată.

De asemenea, deoarece $ f $ este o funcție elementară, avem și continuitatea (pe porțiuni).

Mai trebuie verificată creșterea exponențială, adică faptul că există $ M > 0 $ și $ s_0 \geq 0 $,
cu $ |f(t)| \leq M e^{s_0 t} $, pentru orice $ t \geq 0 $. Această inegalitate este
evidentă, deoarece $ | f(t) | \leq e^t $ pentru $ t $ suficient de mare, deci putem lua
$ s_0 = 1 $, iar $ M $ convenabil ales pentru a satisface condiția.

\vspace{1cm}

2. Calculați transformatele Laplace pentru funcțiile (presupuse original):
\begin{enumerate}[(a)]
\item $ f(t) = 1, t \geq 0 $;
\item $ f(t) = t, t \geq 0 $;
\item $ f(t) = t^n, n \in \NN $;
\item $ f(t) = e^{at}, t \geq 0, a \in \RR $;
\item $ f(t) = \sin(at), t \geq 0, a \in \RR $.
\end{enumerate}

\emph{Soluție:} (a) Avem direct din definiție:
\[
  F(s) = \int_0^\infty e^{-st} dt = \lim_{n \to \infty} \int_0^n e^{-st} dt = \frac{1}{s}, s > 0.
\]

(b) Integrăm prin părți și obținem $ F(s) = \dfrac{1}{s^2} $.

(c) Facem substituția $ st = \tau $ și găsim:
\[
  F(s) = \int_0^\infty e^{-st} t^n dt = \frac{1}{s^{n+1}} \int_0^\infty e^{-\tau} \tau^n d\tau = %
  \frac{n!}{s^{n+1}},
\]
pentru $ s > 0 $, folosind funcția Gamma a lui Euler:
\[
  \Gamma(a) = \int_0^\infty x^{a-1}e^{-x} dx, a > 0, \quad \Gamma(n + 1) = n!, n \in \NN.
\]

(d) $ F(s) = \ds\int_0^\infty e^{at} e^{-st} dt = \ds\int_0^\infty e^{-(s - a)t} dt = %
\dfrac{1}{s-a}, s > a $.

(e) Integrăm prin părți și ajungem la:
\[
  F(s) = \frac{1}{a} - \frac{s^2}{a^2} F(s) \Rightarrow F(s) = \frac{a}{s^2 + a^2}, s > 0.
\]

\vspace{1cm}

3. Folosind tabelul de valori și proprietățile, să se determine transformatele Laplace
pentru funcțiile (presupuse original):
\begin{enumerate}[(a)]
\item $ f(t) = 5 $;
\item $ f(t) = 3t + 6t^2 $;
\item $ f(t) = e^{-3t} $;
\item $ f(t) = 5e^{-3t} $;
\item $ f(t) = \cos (5t) $;
\item $ f(t) = \sin (3t) $;
\item $ f(t) = 3(t-1) + e^{-t-1} $;
\item $ f(t) = 3t^3(t - 1) + e^{-5t} $;
\item $ f(t) = 5e^{-3t} \cos (5t) $;
\item $ f(t) = e^{2t} \sin (3t) $;
\item $ f(t) = te^{-t} \cos (4t) $;
\item $ f(t) = t^2 \sin (3t) $;
\item $ f(t) = t^3 \cos t $.
\end{enumerate}

\emph{Indicații:} În majoritatea cazurilor, se folosește tabelul și proprietatea de liniaritate.
În plus:

(i, j) Folosim $ \Lap (e^{at} f(t)) = F(s - a) $;

(k) Folosim $ \Lap (t f(t)) = -\frac{d}{ds} \Lap(f(t)) $;

(l) Folosim $ \Lap(t^n f(t)) = (-1)^n \dfrac{d^n F(s)}{ds^n} $;

\vspace{1cm}

4. Folosind teorema derivării imaginii, să se determine transformatele Laplace pentru funcțiile
(presupuse original):
\begin{enumerate}[(a)]
\item $ f(t) = t $;
\item $ f(t) = t^2 $;
\item $ f(t) = t\sin t $;
\item $ f(t) = te^t $.
\end{enumerate}

\emph{Indicație:} Conform proprietății de derivare a imaginii, avem:
\[
  \Lap (t f(t)) = -\frac{d}{ds} \Lap (f(t)).
\]

\vspace{1cm}

5. Folosind teorema integrării originalului, să se determine transformatele Laplace pentru
funcțiile (presupuse original):
\begin{enumerate}[(a)]
\item $ f(t) = \ds\int_0^t \cos (2\tau) d\tau $;
\item $ f(t) = \ds\int_0^t e^{3 \tau} \cos(2 \tau) d\tau $;
\item $ f(t) = \ds\int_0^t \tau e^{-3\tau} d\tau $.
\end{enumerate}

\emph{Indicație:} Conform proprietății de integrare a originalului,
avem:
\[
  \Lap \int_0^t f(\tau) d\tau = \frac{F(s)}{s}.
\]

\end{document}
